\introduction

Современный тренировочный процесс всё чаще опирается не только на очное взаимодействие тренера и спортсмена, но и на постоянный обмен цифровыми данными. Тренер должен заранее формировать тренировочные задания, распределять их между спортсменами, учитывать индивидуальные особенности подготовки, получать обратную связь после выполнения тренировки и принимать решения о дальнейшем изменении нагрузки. При работе с одним спортсменом такие действия могут выполняться вручную и не требуют сложной информационной системы. Однако при увеличении количества спортсменов и появлении тренировочных групп ручной обмен сообщениями быстро становится неудобным и ненадёжным.

На практике взаимодействие тренера со спортсменами часто организуется с помощью мессенджеров, электронных таблиц и произвольных файлов. Такой подход не требует внедрения специализированного программного обеспечения, но создаёт ряд существенных проблем, подробно рассмотренных в первой главе. Особенно остро они проявляются в циклических видах спорта, где важно не только получить отдельный отчёт, но и накопить историю структурированных данных для оценки динамики подготовки.

В связи с этим актуальной является задача разработки единой backend-платформы, которая автоматизирует взаимодействие тренера и спортсмена, стандартизирует формат тренировочных отчётов, обеспечивает хранение истории тренировок и предоставляет программные интерфейсы для внешних клиентских приложений. Такая платформа должна решать не только задачу обмена сообщениями, но и задачу организации предметной логики тренировочного процесса: ролей пользователей, связей между тренером и спортсменами, групп, планов, назначений, отчётов, уведомлений и аналитики.

В данной работе рассматривается backend-часть платформы CoachLink. Backend реализует серверную логику, хранение данных, публичные REST API, межсервисное взаимодействие и интеграцию с инфраструктурными компонентами. Мобильное приложение рассматривается только как внешний клиент, который обращается к backend через API Gateway. Разработка мобильного приложения не входит в область ответственности данной выпускной квалификационной работы.

Платформа CoachLink проектируется как расширяемая система для автоматизации тренировочного процесса в различных видах спорта. В целях упрощения разработки прототипа текущая реализация ориентирована исключительно на лёгкую атлетику. При этом архитектура платформы остаётся универсальной: базовые сущности (пользователь, тренер, спортсмен, группа, план, назначение, отчёт, уведомление) не имеют жёсткой привязки к конкретному виду спорта и в дальнейшем могут быть расширены за счёт добавления специфических атрибутов и шаблонов для других дисциплин.

Объектом исследования является процесс информационного взаимодействия между тренером и спортсменом при планировании, выполнении и анализе тренировок.

Предметом исследования является backend-платформа для автоматизации назначения тренировочных заданий, сбора стандартизированных отчётов, отправки уведомлений и анализа данных тренировочного процесса.

Целью выпускной квалификационной работы является разработка открытой микросервисной backend-платформы CoachLink для стандартизации и автоматизации взаимодействия в системе «тренер-спортсмен».

Для достижения поставленной цели необходимо решить следующие задачи:

\begin{enumerate}
\def\labelenumi{\arabic{enumi}.}
\tightlist
\item
  провести анализ предметной области и существующих программных решений, выявить их ограничения в контексте автоматизации взаимодействия тренера и спортсмена;
\item
  сформулировать функциональные и нефункциональные требования к backend-платформе;
\item
  спроектировать микросервисную архитектуру системы, разработать модель данных и спецификации программных интерфейсов;
\item
  реализовать спроектированные микросервисы и обеспечить их взаимодействие;
\item
  провести тестирование backend-компонентов и оценить соответствие разработанной системы исходным требованиям.
\end{enumerate}

При выполнении работы используются методы системного анализа, объектно-ориентированного и компонентного проектирования, проектирования REST API, проектирования реляционных баз данных и тестирования программного обеспечения. Выбор микросервисной архитектуры соотносится с подходами к построению распределённых систем, описанными в работах по шаблонам микросервисов, проектированию независимых сервисов и переходу к cloud-native архитектурам~\cite{richardson-microservices,newman-microservices-2021,dragoni-microservices,balalaie-microservices-devops}.

Теоретическая значимость работы заключается в проектировании модели backend-платформы, объединяющей микросервисный подход, событийное взаимодействие, стандартизированный сбор тренировочных данных и аналитическую обработку этих данных. Работа демонстрирует, как предметная область тренировочного процесса может быть представлена через набор независимых сервисов с чёткими зонами ответственности.

Практическая значимость заключается в создании открытой серверной основы для платформы CoachLink. Реализованный backend может использоваться как база для дальнейшего развития системы: подключения мобильного приложения, добавления новых видов спорта, расширения аналитики, интеграции со спортивными устройствами и развития AI-сценариев. Благодаря открытому API систему можно адаптировать под разные клиентские приложения и внешние интеграции.

Выпускная квалификационная работа состоит из введения, трёх глав, заключения и приложений.

В первой главе проводится аналитический обзор существующих подходов к организации взаимодействия тренера и спортсмена. Рассматриваются мессенджеры, электронные таблицы, спортивные трекеры и специализированные коммерческие платформы. На основе сравнения формулируются требования к backend-платформе.

Во второй главе описывается проектирование CoachLink. Рассматриваются архитектура микросервисов, модель данных, межсервисное взаимодействие, ключевые пользовательские сценарии и выбор технологического стека.

В третьей главе описывается практическая реализация спроектированных микросервисов и результаты тестирования.

В заключении подводятся итоги работы, оценивается достижение цели и формулируются направления дальнейшего развития платформы.
